\documentclass[11pt, a4paper]{article}

\usepackage[utf8]{inputenc}
\usepackage{amsmath}
\usepackage{geometry}
\usepackage{graphicx}
\usepackage{hyperref}

\geometry{a4paper, margin=2cm}

\title{\textbf{Analisi di Missione Interplanetaria: \\ Ottimizzazione delle Finestre di Lancio Terra-Marte tramite Porkchop Plots}}
\author{Il tuo nome}
\date{\today}

\begin{document}

\maketitle

\begin{abstract}
Questo report illustra l'analisi astrodinamica per l'identificazione delle finestre di lancio ottimali in una missione interplanetaria Terra-Marte. Attraverso l'uso delle effemeridi planetarie (JPL) e la risoluzione numerica del Problema di Lambert, è stato sviluppato uno script Python per la generazione di Porkchop Plots, permettendo di minimizzare il $\Delta V$ totale e valutare la fattibilità della missione.
\end{abstract}

\section{Introduzione}
% Contesto delle missioni marziane e importanza delle finestre di lancio.

\section{Inquadramento Teorico}
\subsection{Effemeridi e Posizione Planetaria}
% Spiegazione di come ricaviamo r1 (Terra) e r2 (Marte) usando i dati del JPL.

\subsection{Il Problema di Lambert}
% Formulazione matematica del Boundary Value Problem (BVP).
% Il Teorema di Lambert e il calcolo della traiettoria di trasferimento.

\subsection{Calcolo del $\Delta V$ e Costo Energetico ($C_3$)}
% Come passiamo dalle velocità eliocentriche fornite da Lambert 
% all'effettivo costo di carburante (Delta-V e C3) per il veicolo.

\section{Implementazione Numerica}
% Architettura del codice Python, spiegazione della griglia di partenza/arrivo.

\section{Risultati: Analisi dei Porkchop Plots}
% Qui inseriremo i grafici a curve di livello generati da matplotlib.
% Analizzeremo le famiglie di trasferimento (Type I e Type II).

\section{Conclusioni}

\end{document}